\documentclass[a4paper]{article}
\usepackage{graphicx}
\usepackage{amsmath,amssymb}
\usepackage{hyperref}
\usepackage{minted}
\usepackage{multirow}
\usepackage{float}

\title{Komposition}
\date{}

\begin{document}
    \maketitle
    
    
    \section{Nomen + Nomen \texorpdfstring{$\rightarrow$}{->} Nomen}
        \begin{itemize}
            \item \textbf{Regel:} Zwei Substantive werden zu einem neuen Substantiv zusammengesetzt. Das \textbf{Genus} kommt vom \textbf{letzten} Bestandteil.
            \item \textbf{Beispiele:}
            \begin{itemize}
                \item die Sprache + die Schule $\rightarrow$ die Sprachschule
                \item das Buch + der Laden $\rightarrow$ der Buchladen
                \item die Arbeit + das Zimmer $\rightarrow$ das Arbeitszimmer
            \end{itemize}
        \end{itemize}
    
    
    \section{Adjektiv + Nomen \texorpdfstring{$\rightarrow$}{->} Nomen}
        \begin{itemize}
            \item \textbf{Regel:} Ein Adjektiv bestimmt ein Substantiv, zusammen entsteht ein neues Substantiv.
            \item \textbf{Beispiele:}
            \begin{itemize}
                \item hoch + die Schule $\rightarrow$ die Hochschule
                \item klein + das Kind $\rightarrow$ das Kleinkind
            \end{itemize}
        \end{itemize}
    
    
    \section{Adjektiv + Adjektiv \texorpdfstring{$\rightarrow$}{->} Adjektiv}
        \begin{itemize}
            \item \textbf{Regel:} Zwei Adjektive werden zu einem neuen Adjektiv zusammengesetzt.
            \item \textbf{Beispiele:}
            \begin{itemize}
                \item dunkel + blau $\rightarrow$ dunkelblau
                \item hell + grün $\rightarrow$ hellgrün
            \end{itemize}
        \end{itemize}
    
    
    \section{Verb + Verb \texorpdfstring{$\rightarrow$}{->} Verb}
        \begin{itemize}
            \item \textbf{Regel:} Zwei Verben bilden zusammen ein neues Verb; häufig ist das zweite Verb \textit{lernen}, \textit{lassen} oder \textit{gehen}.
            \item \textbf{Beispiele:}
            \begin{itemize}
                \item kennen + lernen $\rightarrow$ kennenlernen
                \item sitzen + bleiben $\rightarrow$ sitzenbleiben
            \end{itemize}
        \end{itemize}
    
    
    \section{Nomen + Adjektiv \texorpdfstring{$\rightarrow$}{->} Adjektiv}
        \begin{itemize}
            \item \textbf{Regel:} Ein Substantiv (oft ohne Endung) verbindet sich mit einem Adjektiv und bildet ein neues Adjektiv.
            \item \textbf{Beispiele:}
            \begin{itemize}
                \item die Welt + bekannt $\rightarrow$ weltbekannt
                \item der Alltag + tauglich $\rightarrow$ alltagstauglich
            \end{itemize}
        \end{itemize}
    
    
    \section{Fugenelemente in der Komposition}
        \begin{itemize}
            \item \textbf{Regel:} Zwischen den Teilen steht manchmal ein \textbf{Fugenelement} (z.\,B. \textit{-s-}, \textit{-n-}, \textit{-en-}, \textit{-es-}). Es erleichtert oft die Aussprache oder ist historisch bedingt.
            \item \textbf{Beispiele:}
            \begin{itemize}
                \item die Liebe + das Lied $\rightarrow$ das Liebeslied
                \item die Arbeit + das Zimmer $\rightarrow$ das Arbeitszimmer
                \item der Name + das Schild $\rightarrow$ das Namensschild
            \end{itemize}
        \end{itemize}
    % \bibliographystyle{plain}
    % \bibliography{/home/donkarlo/Dropbox/repo/shared_research/bibliography/bibliography.bib}
\end{document}